\documentclass[a4paper,11pt]{article}
\usepackage{exam-style-842-unit}
\examinehead{842 信息化概论 · 第 3 单元练习册}{练习二 参考答案与解析}
\begin{document}

\unitanswerscover{842 信息化概论 · 单元练习册}{第 3 单元参考答案}{练习二（通信与网络技术）的答案与解析；每道判断题均给出改正后的正确表述。}

\answerpart{一、单项选择题（10 题，每题 3 分，共 30 分）}

\q B。\begin{analysis}计算机网络的发展大致经历四个阶段：诞生阶段（20 世纪 60 年代中期以前，以单个计算机为中心的远程联机系统）、形成阶段（60 年代中期至 70 年代，多个主机互联，典型代表是 ARPANET）、互联互通阶段（70 年代末至 90 年代，统一体系结构并遵守国际标准，出现 TCP/IP 与 OSI）、高速网络技术阶段（90 年代至今，以 Internet 为代表）（教材 6.2.1，第 102—103 页）。A 诞生阶段的典型应用是飞机订票系统；C 互联互通阶段的标志是 TCP/IP 与 OSI 体系结构；D 高速网络技术阶段的代表是 Internet。\end{analysis}

\q B。\begin{analysis}三网融合是指电信网、广播电视网、互联网在向宽带通信网、数字电视网、下一代互联网演进过程中，技术功能趋于一致、业务范围趋于相同，网络互联互通、资源共享，能为用户提供语音、数据和广播电视等多种服务；它不意味着三大网络的物理合一，而主要是高层业务应用的融合（教材 6.2.2，第 103 页）。A 与教材表述相反；C 是企业经营行为；D 终端融合只是三网融合的一个层面。\end{analysis}

\q A。\begin{analysis}网络协议是为进行计算机网络中的数据交换而建立的规则、标准或约定的集合，其三个要素是语义（涉及协调与差错处理的控制信息）、语法（涉及数据及控制信息的格式、编码及信号电平等）和同步（涉及速度匹配和排序等）（教材 6.3.1，第 105 页）。B、C 是信息系统的组成与功能；D 是通信过程的具体动作，不是协议要素。\end{analysis}

\q C。\begin{analysis}网络体系结构就是计算机网络各层次及其协议的集合（教材 6.3.1，第 105 页）。A 网络拓扑描述网络的物理连接形状；B 网络协议只是某一层的规则集合；D 网络操作系统是管理网络资源的系统软件，均不是“层次及其协议集合”的定义。\end{analysis}

\q B。\begin{analysis}运输层是第一个端—端、即主机—主机的层次，它提供端到端的透明数据运输服务，使高层用户不必关心通信子网的存在，并处理端到端的差错控制和流量控制（教材 6.3.1，第 107—108 页）。A 网络层关心通信子网的运行控制与路由选择；C 会话层是进程—进程的层次；D 表示层处理语法表示变换。\end{analysis}

\q B。\begin{analysis}表示层为上层用户提供共同的数据或信息的语法表示变换，负责把计算机内部的表示形式转换成网络通信中采用的标准表示形式，数据压缩和加密也是表示层可提供的表示变换功能（教材 6.3.1，第 108 页）。A 会话层负责对话连接与同步；C 应用层面向具体网络应用；D 数据链路层负责帧的传输与差错控制。\end{analysis}

\q A。\begin{analysis}TCP/IP 的网络层（互联层）主要实现网际互联，将数据包封装成 Internet 数据报并运行必要的路由算法，其中网际协议 IP 负责在主机和网络之间寻址和路由数据包（教材 6.3.1，第 108 页）。B TCP 和 C UDP 是传输层协议；D HTTP 是应用层协议，均不负责寻址和路由。\end{analysis}

\q B。\begin{analysis}传输层主要有两个协议：传输控制协议 TCP 为应用程序提供可靠的通信连接，适合一次传输大批数据并适用于要求得到响应的应用；用户数据报协议 UDP 提供无连接通信，不对传送包进行可靠保证，适合一次传输小量数据，可靠性由应用层负责（教材 6.3.1，第 108 页）。A、C、D 与题干描述不符。\end{analysis}

\q A。\begin{analysis}SNMP 是基于 TCP/IP 协议族的网络管理标准，由一组网络管理的标准组成，包含一个应用层协议、数据库模型和一组资源对象（教材 6.3.2，第 109—110 页）；在 TCP/IP 协议族图中，SNMP 与 FTP、HTTP、SMTP、DNS 等一起列在应用层。B 传输层、C 网络层、D 网络接口层均不是 SNMP 所属层次。\end{analysis}

\q A。\begin{analysis}网络安全具有三个基本属性：机密性、完整性和可用性。机密性指保证信息与信息系统不被非授权者获取与使用，主要措施是密码技术；完整性指保证信息不被篡改或破坏；可用性指保证信息与信息系统可被授权人正常使用（教材 6.3.3，第 110—111 页）。B 是信息共享层面的表述；C 是性能与经济性要求；D 是系统结构特征，均不是安全属性。\end{analysis}

\answerpart{二、多项选择题（3 题，每题 4 分，共 12 分）}

\q ABCD。\begin{analysis}三网融合在概念上从不同角度和层次分析，可以涉及技术融合、业务融合、行业融合、终端融合及网络融合（教材 6.2.2，第 103—104 页），四项全对。\end{analysis}

\q ABC。\begin{analysis}OSI 七层模型从下到上为物理层、数据链路层、网络层、运输层、会话层、表示层和应用层（教材 6.3.1，第 106 页）。A、B、C 都在其中；D 错误：网络接口层是 TCP/IP 四层模型中的层次（又称数据链路层），不是 OSI 的层次名称。\end{analysis}

\q ABC。\begin{analysis}网络安全技术机制主要包括：加密和隐藏、认证和授权、审计和定位、完整性保证、权限和存取控制、任务填充（教材 6.3.3，第 111 页）。D 错误：“随意开放共享”与安全机制背道而驰，访问必须经过认证授权与存取控制，不能随意开放。\end{analysis}

\answerpart{三、判断题（6 题，每题 2 分，共 12 分）}

\q ×。改正：三网融合并不意味着三大网络的物理合一，而主要是指高层业务应用的融合。\begin{analysis}教材 6.2.2（第 103 页）明确指出：三网融合并不意味着三大网络的物理合一，而主要是指高层业务应用的融合；三大网络通过技术改造实现技术功能趋于一致、业务范围趋于相同、互联互通、资源共享。\end{analysis}

\q √。\begin{analysis}现代移动通信已经历四个阶段：第一代（1G）是模拟、仅限语音的蜂窝电话标准，制定于 20 世纪 80 年代，容量有限、制式太多、互不兼容、保密性差、通话质量不高，不能提供数据业务和自动漫游；第二代（2G）以数字语音传输技术为核心；第三代（3G）以宽带 CDMA 技术为主，能同时提供语音和数据业务；第四代（4G）集 3G 与 WLAN 于一体，包括 TD-LTE 和 FDD-LTE 两种制式（教材 6.1.2，第 98 页）。\end{analysis}

\q ×。改正：OSI 参考模型并非具体实现的描述，它只是一个为制定标准而提供的概念性框架，只有各种协议是可以实现的。\begin{analysis}教材 6.3.1（第 106 页）强调：OSI 参考模型并非具体实现的描述，它只是一个为制定标准而提供的概念性框架；在 OSI 中只有各种协议是可以实现的，网络设备只有与 OSI 及有关协议一致时才能互连。\end{analysis}

\q √。\begin{analysis}运输层是第一个端—端、即主机—主机的层次，提供端到端的透明数据运输服务，使高层用户不必关心通信子网的存在（教材 6.3.1，第 107—108 页）。\end{analysis}

\q ×。改正：TCP 是面向连接的可靠传输协议，UDP 是无连接、不对传送包作可靠保证的传输协议。\begin{analysis}教材 6.3.1（第 108 页）：TCP 为应用程序提供可靠的通信连接，适合一次传输大批数据；UDP 提供无连接通信，且不对传送包进行可靠的保证，可靠性由应用层负责。题干把二者说反了。\end{analysis}

\q √。\begin{analysis}网络安全具有三个基本属性——机密性、完整性和可用性（教材 6.3.3，第 110—111 页）。广义上网络安全还涉及可控性，但三个基本属性是机密性、完整性和可用性。\end{analysis}

\answerpart{四、配对题（6 对，每对 1 分，共 6 分）}

\q （1）—A\quad （2）—B\quad （3）—C\quad （4）—D\quad （5）—E\quad （6）—F。
\begin{analysis}OSI 各层主要功能：物理层使原始比特流在物理媒体上传输；数据链路层把比特流组织成帧，通过校验、确认和反馈重发把不可靠物理链路改造成对网络层无差错的数据链路；网络层以分组为单位传输，解决路由选择、拥塞控制和网际互连；运输层提供端到端的透明数据传输；应用层为特定类型的网络应用提供访问 OSI 环境的手段；SNMP 是基于 TCP/IP 协议族的网络管理标准。（教材 6.3.1、6.3.2，第 105—110 页）\end{analysis}

\answerpart{五、简答题（4 题，每题 6 分，共 24 分）}

\q 参考答案要点：
（1）含义：三网融合是指电信网、广播电视网、互联网在向宽带通信网、数字电视网、下一代互联网演进过程中，三大网络通过技术改造，其技术功能趋于一致，业务范围趋于相同，网络互联互通、资源共享，能为用户提供语音、数据和广播电视等多种服务（教材 6.2.2，第 103 页）。
（2）实质：三网融合并不意味着三大网络的物理合一，而主要是指高层业务应用的融合；在手机上可以看电视、上网，用电视可以打电话、上网，用计算机也可以打电话、看电视，三者相互交叉、形成你中有我、我中有你的格局。
（3）融合层面：涉及技术融合、业务融合、行业融合、终端融合及网络融合。
（4）支撑技术：数字技术（各种业务统一编码为比特流）、宽带技术（以光纤通信为主体，提供大容量传输平台）、软件技术（使三大网络及其终端通过软件变更支持各种业务）、IP 技术（特别是 IPv6，在内容与传送介质之间建立统一的映射）。
\begin{analysis}得分点：含义 2 分，实质（不等于物理合一）1 分，五个层面 1 分，四项支撑技术 2 分。常见失分：把三网融合理解为三网物理合并；漏答“高层业务应用融合”这一实质；只写技术融合而漏掉业务、行业、终端、网络融合。教材 6.2.2，第 103—104 页。\end{analysis}

\q 参考答案要点：
OSI 参考模型从下到上分为七层（教材 6.3.1，第 106—109 页）：
（1）物理层：定义建立、维护和拆除物理链路所需的机械的、电气的、功能的、规程的特性，使原始比特流能在物理媒体上传输。
（2）数据链路层：把比特流组织成帧，通过校验、确认和反馈重发把不可靠物理链路改造成对网络层无差错的数据链路，并进行流量控制。
（3）网络层：以分组为单位传输，解决数据分组跨越通信子网从源到目的地的路由选择、拥塞控制和网际互连问题。
（4）运输层：第一个端到端（主机—主机）层次，提供端到端的透明数据运输服务，处理端到端的差错控制和流量控制。
（5）会话层：进程—进程的层次，组织和同步不同主机上各种进程间的通信（对话），建立和拆除对话连接，提供同步点机制。
（6）表示层：为上层用户提供数据或信息的语法表示变换，并具有数据压缩和加密功能。
（7）应用层：开放系统互连环境的最高层，为特定类型的网络应用提供访问 OSI 环境的手段，如文件传送访问和管理、文电处理系统、虚拟终端等。
\begin{analysis}得分点：七层名称与顺序 3 分，各层功能 3 分。常见失分：层序颠倒（把表示层放到会话层之下）；把运输层说成“网络层的一部分”；各层功能张冠李戴（如把路由选择说成数据链路层功能）。教材 6.3.1，第 105—109 页。\end{analysis}

\q 参考答案要点：
（1）网络管理的概念：网络管理包括对硬件、软件和人力的使用、综合与协调，以便对网络资源进行监视、测试、配置、分析、评价和控制，以合理的价格满足网络的实时运行性能、服务质量等需求；当网络出现故障时能及时报告和处理，并协调、保持网络系统的高效运行。网络管理是对网络运行状态进行监测和控制，监测是控制的前提，控制是监测的结果（教材 6.3.2，第 109 页）。
（2）网络管理系统的结构：通常分为管理站和代理系统两部分。
（3）网络管理的主要功能：
（1）网络监视：收集系统和子网的状态信息，分析被管理设备的行为。包括性能监视（可用性、响应时间、正确性、吞吐率、利用率）、故障监视（故障检测、报警、预测、诊断和定位）、计费监视（单位时间内对资源的占用和访问）。
（2）网络控制：修改设备参数或重新配置网络资源以改善运行状态，包括配置控制（初始化、维护和关闭设备，定义和修改互联关系等）和安全控制（对数据、硬件、软件、信息、用户访问的安全控制）。
（4）SNMP 的组成：SNMP 管理的网络主要由被管理的设备、SNMP 代理和网络管理系统（NMS）三部分组成；网络中每一个被管理设备都存在一个管理信息库（MIB）用于收集并储存管理信息，NMS 通过 SNMP 获取这些信息（教材 6.3.2，第 110 页）。
\begin{analysis}得分点：网络管理概念 2 分，管理站与代理结构 1 分，监视与控制功能 2 分，SNMP 三部分组成与 MIB 1 分。常见失分：只答“用软件监控网络”而无监测与控制之分；漏答管理信息库（MIB）；把 SNMP 与管理站混为一谈。教材 6.3.2，第 109—110 页。\end{analysis}

\q 参考答案要点：
（1）网络安全的概念：网络安全从本质上讲就是网络上信息的安全；广义上讲，凡涉及网络上信息的保密性、完整性、可用性和可控性的相关技术和理论都是网络安全的研究领域。网络安全是指网络系统的硬件、软件及数据受到保护，不遭受偶然或恶意的破坏、更改、泄露，系统连续可靠正常地运行，网络服务不中断（教材 6.3.3，第 110—111 页）。
（2）三个基本属性：机密性（信息与系统不被非授权者获取与使用，主要靠密码技术）、完整性（信息不被篡改或破坏）、可用性（信息与系统可被授权者正常使用）。
（3）网络安全涉及的内容：运行系统安全、网络上系统信息的安全、网络上信息传输的安全、网络上信息内容的安全。
（4）主要技术机制：加密和隐藏、认证和授权、审计和定位、完整性保证、权限和存取控制、任务填充（教材 6.3.3，第 111 页）。
（5）网络安全不仅是技术问题，更是管理问题：除技术防控外，还必须制订网络运营过程中的安全管理措施，确定正确的网络运营目标，依据相关法律设计可行的网络管理方案和网络信息安全管理制度。
\begin{analysis}得分点：概念 2 分，三属性 2 分，技术机制答出四项以上 1.5 分，指出管理层面 0.5 分。常见失分：只答“防病毒、防火墙”；把机密性、完整性、可用性写错或漏项；忽视管理与法律层面。教材 6.3.3，第 110—112 页；本题对应 2021 年真题“网络安全机制的主要内容”。\end{analysis}

\answerpart{六、综合题（论述，1 题，共 16 分）}

\q 参考答案要点：
（1）数据传输过程（4 分）：园区传感器采集温湿度、光照等数据后，数据经接入网送到园区网络设备，按 TCP/IP 四层模型自上而下封装：应用层由数据采集与监控程序按 HTTP、MQTT 等应用协议组织数据；传输层用 TCP 或 UDP 加上端口号，实现端到端通信（监控指令等可靠性要求高的用 TCP，实时视频可用 UDP）；网络层由 IP 协议加上源和目的 IP 地址，封装成数据报并进行路由选择，经园区核心交换机、路由器或网关转发；网络接口层（数据链路层）把数据封装成帧，通过以太网、光纤或无线链路发送。数据到达用户终端后，再按相反顺序自下而上逐层解封装，交应用层显示或处理。这一过程体现了 OSI/TCP/IP 分层、对等层通信和逐层封装的思想。
（2）三网融合的体现（4 分）：园区用同一光纤网络同时承载互联网数据、有线电视信号和 IP 电话，技术功能趋于一致、业务范围趋于相同，实现网络互联互通、资源共享，用户可以通过手机或电脑看电视、上网，通过电视或计算机打电话，体现了技术融合、业务融合和终端融合；但三大网络并未合并为一张物理网络，融合主要体现在高层业务应用层面。数字技术、宽带技术、软件技术和 IP 技术是三网融合得以实现的支撑。
（3）网络管理（4 分）：在园区部署网络管理站，在各交换机、路由器、服务器上运行 SNMP 代理，建立管理信息库（MIB），对设备进行监视与控制。网络监视包括性能监视（带宽利用率、时延、丢包率、设备 CPU 与内存占用）、故障监视（设备告警、链路中断的检测、定位与诊断）和计费监视；网络控制包括配置控制（VLAN 划分、端口与路由配置、设备参数修改）和安全控制（访问控制、账号权限管理）。同时建立运维值班、巡检、变更和应急响应制度，使网络管理向综合化、自动化、智能化发展。
（4）网络安全措施（4 分）：以机密性、完整性、可用性为目标。技术措施：数据传输和存储采用加密技术，敏感数据脱敏；实行身份认证与授权、基于角色的访问控制和最小权限原则；利用数字签名、校验码保证数据完整性；部署防火墙、入侵检测与防病毒系统；建立数据备份与灾难恢复机制；通过日志审计和安全监测发现异常访问。管理措施：制定园区网络安全管理制度和应急预案，明确数据分级分类与共享审批流程，落实安全责任制，加强人员安全培训，遵守网络安全与数据保护相关法律法规。只有技术与管理并重，才能保证园区网络连续可靠运行、数据不被泄露和篡改。
\begin{analysis}得分点：（1）分层封装与解封装过程 4 分；（2）三网融合含义与在本案例中的体现（同一网络承载多业务、高层业务融合、不等于物理合一）4 分；（3）管理站与代理、SNMP/MIB、监视与控制功能 4 分；（4）安全三属性＋技术机制＋管理制度 4 分。常见失分：把“三网融合”答成“三网合一”；网络管理只答“派人维护”；安全部分只列技术手段而漏掉管理制度与法律依据。教材 6.2.2、6.3.1、6.3.2、6.3.3（第 103—112 页）。\end{analysis}

\end{document}
